\documentclass[11pt,a4paper]{article}

\usepackage[utf8]{inputenc}
\usepackage[T1]{fontenc}
\usepackage[french]{babel}
\usepackage{amsmath}
\usepackage[margin=2.3cm]{geometry}
\usepackage{booktabs}
\usepackage{array}
\usepackage{ragged2e}
\usepackage{titlesec}
\usepackage{setspace}
\usepackage[colorlinks=true,linkcolor=black,urlcolor=black]{hyperref}

\titleformat{\section}{\large\bfseries}{\thesection. }{0pt}{}
\setstretch{1.12}
\newcolumntype{L}[1]{>{\RaggedRight\arraybackslash}p{#1}}
\renewcommand{\arraystretch}{1.25}

\begin{document}

% ======================= TITRE =======================
\begin{center}
  {\LARGE\bfseries POURQUOI 291 INTERVIEWS SONT INCOMPLETES\par}
  \vspace{0.3em}
  {\large et comment les rattraper\par}
  \vspace{0.4em}
  {\normalsize Note de diagnostic --- analyse de la completude logique COSO\par}
\end{center}
\vspace{0.5em}
\hrule
\vspace{0.8em}

Sur les \textbf{1032 interviews} de la base finale : 627 completes,
62 quasi-completes, \textbf{10 moyennes}, \textbf{291 incompletes} et
42 vides. Cette note explique pourquoi les interviews des classes
\textbf{4-INCOMPLET} (291) et \textbf{3-MOYEN} (10) sont incompletes, et
quelles actions de rattrapage sont possibles.

% ======================= 1. LES 291 =======================
\section{Pourquoi les 291 interviews « 4-INCOMPLET » le sont}

Elles n'ont repondu qu'a \textbf{1 a 29 questions} sur $\sim$65 attendues
(moyenne 3,2). La progression section par section est sans ambiguite
(tableau~\ref{tab:prog291}).

\begin{table}[h]
\centering
\caption{Progression des 291 interviews incompletes par section.}
\label{tab:prog291}
\small
\begin{tabular}{L{3.4cm}L{3.4cm}c}
\toprule
\textbf{Section} & \textbf{1re question} & \textbf{Interviews ayant repondu} \\
\midrule
A (contact)          & A1   & 97,3~\% \\
B (identification)   & B1   & 100,0~\% \\
C                    & C1   & 2,7~\%  \\
D                    & D1   & 0,7~\%  \\
E                    & E1   & 0,3~\%  \\
EMPLOYE              & EMPLOYE & 0,7~\% \\
F                    & F1   & 0,0~\%  \\
G                    & G1   & 0,0~\%  \\
H                    & H1   & 0,0~\%  \\
J                    & J1   & 0,0~\%  \\
K                    & K1   & 0,0~\%  \\
L                    & L1   & 0,0~\%  \\
M                    & M1   & 0,0~\%  \\
N                    & N1   & 0,0~\%  \\
O                    & O1   & 0,0~\%  \\
P                    & P1   & 0,0~\%  \\
Q                    & Q1   & 0,0~\%  \\
R                    & R1   & 0,0~\%  \\
\bottomrule
\end{tabular}
\end{table}

\noindent Lecture : les sections A (contact) et B (identification) sont
remplies a presque 100~\%, puis tout s'arrete des le debut de la section
C : seules \textbf{2,7~\%} des interviews ont une reponse a C1 (8 sur 291).
Ce sont donc des \emph{abandons precoces} : l'interview a ete ouverte puis
recloturee sans etre menee, et non des abandons en cours de questionnaire.

La repartition par statut (tableau~\ref{tab:statut291}) montre que cela
touche toutes les categories, y compris \textbf{78 interviews pourtant
classees « validees »} (statut 100).

\begin{table}[h]
\centering
\caption{Questions remplies dans les 291 interviews incompletes, par statut.}
\label{tab:statut291}
\small
\begin{tabular}{ccccc}
\toprule
\textbf{Statut} & \textbf{Nb interviews} & \textbf{min} & \textbf{moyenne} & \textbf{max} \\
\midrule
60 & 99  & 1 & 3,2 & 29 \\
65 & 114 & 1 & 2,4 & 5  \\
100 & 78 & 1 & 4,3 & 27 \\
\bottomrule
\end{tabular}
\end{table}

\noindent Concernant le consentement : sur les 291, \textbf{107 ont consenti},
29 ont refuse et 155 n'ont pas de reponse de consentement enregistree. Sans
consentement, aucun recontact legal ; avec consentement, un rappel est
theoriquement possible.

% ======================= 2. LES 10 =======================
\section{Pourquoi les 10 interviews « 3-MOYEN » le sont}

Profil inverse : elles sont allees loin (de 43 a 90 questions remplies sur
$\sim$90 attendues, moyenne 69) mais ont abandonne sur les sections de fin
(tableau~\ref{tab:prog10}).

\begin{table}[h]
\centering
\caption{Progression des 10 interviews moyennes par section.}
\label{tab:prog10}
\small
\begin{tabular}{L{3.4cm}L{3.4cm}c}
\toprule
\textbf{Section} & \textbf{1re question} & \textbf{Interviews ayant repondu} \\
\midrule
A, B, C, D, E, EMPLOYE, H, K & -- & 100,0~\% \\
L & L1 & 90,0~\% \\
M & M1 & 70,0~\% \\
N & N1 & 70,0~\% \\
F, G, I & F1, G1 & 60,0~\% \\
J & J1 & 30,0~\% \\
O & O1 & 50,0~\% \\
Q & Q1 & 50,0~\% \\
R & R1 & 40,0~\% \\
P & P1 & 0,0~\% \\
\bottomrule
\end{tabular}
\end{table}

\noindent Lecture : sections A a E, EMPLOYE, H et K remplies a 100~\% ; les
sections de la fin (O, P, Q, R notamment) ne sont remplies qu'a 40--60~\%.
Bilan : ce sont de vrais abandons en fin de questionnaire, sur des blocs
conditionnels et les observations de l'enqueteur.

% ======================= 3. RATTRAPAGE =======================
\section{Comment rattraper (selon la classe)}

\begin{table}[h]
\centering
\caption{Actions de rattrapage par classe de completude.}
\label{tab:rattrapage}
\small
\begin{tabular}{L{4.2cm}L{4.6cm}L{5.2cm}}
\toprule
\textbf{Classe} & \textbf{Etat} & \textbf{Action de rattrapage} \\
\midrule
2-QUASI\_COMPLET (62)   & il manque 1 a 3 questions &
  recontact cible : relancer precise-ment les questions manquantes \\
3-MOYEN (10)            & fin de questionnaire non passee &
  recontact cible de la fin (sections O a R) --- petit volume, realisable \\
4-INCOMPLET (291)       & reponses presque inexistantes &
  recontact complet si consentement (107) sinon exclusion de l'analyse \\
5-VIDE (42)             & aucun champ exploitable &
  exclusion de l'analyse \\
\bottomrule
\end{tabular}
\end{table}

\noindent Ouverture pratique : pour les interviews a relancer, la liste
exacte des identifiants et des questions manquantes est dans le fichier
\texttt{resultats/manquants\_id\_detail.xlsx} (filtrer par classe puis par
\texttt{interview\_\_id}).

\end{document}